% 第九卷 14.2.2.2：AlphaCode 的大规模采样、示例测试过滤与按行为聚类（跨栏，约160 mm宽）。
% FigureInput：题面与测试输入；FigureModel：生成模型；FigureOutput：提交集合；中间为执行与筛选步骤。
\begin{tikzpicture}[
  x=1mm,y=1mm,
  figure picture,
  srcbox/.style={figure node,draw=FigureInput,fill=FigureInputFill,minimum height=15mm,
    minimum width=20mm,inner xsep=2pt},
  genbox/.style={figure node,draw=FigureModel,fill=FigureModelFill,minimum height=15mm,
    minimum width=29mm,inner xsep=2pt,line width=.9pt},
  procbox/.style={figure node,draw=FigureStroke,fill=white,minimum height=15mm,
    minimum width=29mm,inner xsep=2pt},
  subbox/.style={figure node,draw=FigureOutput,fill=FigureOutputFill,minimum height=15mm,
    minimum width=24mm,inner xsep=2pt}
]
  \node[srcbox] (prob) at (10,24) {题面\\示例输入输出};
  \node[genbox] (gen) at (43,24) {代码生成模型\\采样$n$个程序\\语言、标签、难度随机};
  \node[procbox] (filt) at (79,24) {示例测试过滤\\约99\%被淘汰};
  \node[procbox] (clus) at (115,24) {在新输入上执行\\按输出向量聚类};
  \node[subbox] (sub) at (150,24) {由大到小\\每类取一个\\至多$K$次提交};
  \node[srcbox,minimum width=29mm] (tin) at (115,3) {测试输入生成模型\\产生$m$个新输入};

  \draw[figure flow] (prob) -- (gen);
  \draw[figure flow] (gen) -- node[figure edge label,above=1pt] {$n$个} (filt);
  \draw[figure flow] (filt) -- node[figure edge label,above=1pt] {$n'$个} (clus);
  \draw[figure flow] (clus) -- node[figure edge label,above=1pt] {$C$类} (sub);
  \draw[figure flow] (tin) -- (clus);
  \draw[figure flow,dash pattern=on 2pt off 1.5pt] (prob.south) |- node[figure note,pos=.75,above=1pt] {以题面为条件} (tin.west);
\end{tikzpicture}
